\section{Exact long matching and the secondary scale}
\label{rw:secondary-scale}

This section treats the remaining corner
\[
 \bar B>0,\qquad \bar c=\bar m=0,\qquad
 |\bar a|\le2\sqrt{\bar B},\qquad p=q=\infty.
\]
The primary affine upper field is parabolic; its returning arcs may
have arbitrarily long passages on a scale determined by $\bar d$ and
$\bar m$, independently of the primary scale. Follow each arc to its
first exit, at either a fixed-height wedge section or a secondary-scale
section. The former gives two long
legs joined by a regular vertical passage; the latter gives a normalized
affine upper problem. Attach the original lower passages in both cases
and compare their outputs at the same central section.
Throughout, $\bar B$ belongs to a fixed compact positive interval.
Write
\[
 Z=Y+\epsilon,\quad \bar d=\bar c+\nu,\quad
 \theta=\epsilon\bar B,\quad b=1-\theta,\quad k=1-2\theta,\quad
 \beta_\epsilon=k/b,\quad h_I=b/\bar B,\quad \mu=b/\theta.
\]
In particular $\beta_\epsilon\mu=k/\theta$ exactly.
The secondary scale is introduced at the second exit. Every counting
interval is restricted by the complete lower passing conditions and
positive original central heights. Coefficient division uses the signed
central continuations with reference values fixed independently of input.

\subsection{The weighted field and its exact long passage}

The original upper equation, with its forcing retained, is
\begin{equation}\label{rw:secondary:retained}
 Z_\xi=Z[X+\bar d(Z-\epsilon)]-\epsilon\nu(Z-\epsilon).
\end{equation}
This identity is the reason for using $\bar d$, rather than $\bar c$,
in the new scale. For $r=Z^{-1/2}$, $v=Xr$ and
$\dd\tau/\dd\xi=1/r$, put
\[
 D_v=\bar d/r,\quad M_v=\bar m/r^2,\quad
 A_v=\bar a r,\quad N_v=\nu r,\quad E_v=\epsilon r^2.
\]
The exact polynomial equations are
\begin{equation}\label{rw:secondary:vertical}
\begin{gathered}
 H_v=v+D_v(1-E_v)-\epsilon N_v(1-E_v),\\
 P_v=(1-E_v)[M_v(1-E_v)-1]+\theta v^2+\epsilon A_vv-vH_v/2,\\
 r_\tau=-rH_v/2,\quad v_\tau=P_v,\quad
 (D_v)_\tau=D_vH_v/2,\quad (M_v)_\tau=M_vH_v,\\
 (A_v)_\tau=-A_vH_v/2,\quad (N_v)_\tau=-N_vH_v/2,\quad
 (E_v)_\tau=-E_vH_v.
\end{gathered}
\end{equation}
Conserved products recover the original coefficients, for example
$D_vr=\bar d$ and $N_v/r=\nu$, so retain these variables in the
vertical passage.

For any fixed $V>0$, small extended entries at $v=V$ pass uniquely
to $v=-V$, with finite time and bounded finite jets. Indeed, at zero
extended variables $P_v=-1-kv^2/2$. On a slightly larger compact tube keep
$P_v<-1/2$, bound the logarithmic rates of the five extended variables,
and reduce the entry widths by their growth factor for time $4V$.
First exit then gives passage in time at most $4V$.
Finite variational equations and the transverse hitting equation give
the stated jets, including the inverse. The radial component is
$r_-=r_+L_v$ with positive bounded unit
$L_v=\exp[-\frac12\int H_v\,\dd\tau]$.
Reflection makes $L_v=1$ on the reversible leaves.
These are jets in extended entries, not bounded original
$\bar d$ derivatives at $r_+=0$.

On the outgoing adjoining chart set
\[
 \varphi=X/Z,\quad R=1/Z=\varphi^2w,\quad
 g=\bar d/\varphi,\quad j=\bar m/\varphi^2,\quad
 e=\epsilon\varphi^2w .
\]
Divide the field by $X>0$. Direct substitution gives
\begin{equation}\label{rw:secondary:adjoining}
\begin{gathered}
 G=1+g(1-e)-\epsilon\nu\varphi w(1-e),\\
 W=(1-e)[j(1-e)-w]-b-g(1-e)
                     +\epsilon\varphi w[\bar a+\nu(1-e)],\\
 \varphi'=\varphi W,\qquad w'=-w(G+2W),\qquad
 g'=-gW,\qquad j'=-2jW.
\end{gathered}
\end{equation}
Thus $g\varphi=\bar d$, $j\varphi^2=\bar m$ remain constant.
Its origin has contracting and expanding rates $-b,k$.

\begin{proposition}[The exact long map]\label{rw:exact-long}
Fix an order $N$, a section $\varphi_0>0$, and a finite terminal
height range $0\le w_e\le w_*$. On sufficiently small common
ordinary coefficient neighborhoods, the backward map from
$(\varphi,w,g,j)=(H,w_e,\gamma,J)$ to $\varphi=\varphi_0$ is
\begin{equation}\label{rw:secondary:long-map}
 w_{\rm in}=w_e(H/\varphi_0)^{\beta_\epsilon}e^\Psi .
\end{equation}
For one $q>0$ depending only on the fixed order,
\begin{equation}\label{rw:secondary:long-jets}
\begin{aligned}
 \Psi&=\Psi_0(\gamma,J,w_e;\theta)+O_N^{\rm aux}(H^q),\qquad
 H\partial_H\Psi=O_{N-1}^{\rm aux}(H^q),\\
 \partial_x\Psi&=\epsilon S_x,\qquad
 S_x=O_{N-1}^{\rm aux}(H^q),\qquad x\in\{\bar a,\nu\}.
\end{aligned}
\end{equation}
In the auxiliary norms $\theta$ and $\epsilon$ are independent.
The derivatives are ordinary terminal-height derivatives through zero,
ordinary $\gamma,J,\bar a,\nu,\theta,\epsilon$ derivatives and
$H\partial_H$, with the indicated total-order budgets.
Ordinary $\epsilon$ differentiation of $\epsilon S_x$ may remove its
explicit factor. After $\theta=\epsilon\bar B$ is imposed, scale
derivatives are taken with
$E_\epsilon=\epsilon\partial_\epsilon+\theta\partial_\theta$;
they preserve that factor. No bound on $\epsilon|\log H|$ is required.
\end{proposition}

\begin{proof}
Put $t=\log(\varphi/H)$ and $L=\log(\varphi_0/H)$. Exactly,
\begin{equation}\label{rw:secondary:long-scalar}
 \varphi=He^t,\quad g=\gamma e^{-t},\quad j=Je^{-2t},\qquad
 w_t=-wQ,\quad Q=2+G/W,\quad 0\le t\le L .
\end{equation}
Choose the larger rectangle with $W<-c$, $Q>1/2$; then
$w\le w_e e^{-t/2}$ and no earlier exit is possible.
Subtracting the exact saddle ratio gives
\[
 Q-\beta_\epsilon=(G-1)/W+(W+b)/(bW)
\]
with a right side integrable uniformly in $L$. Integration gives
$\Psi=-\int_0^L(Q-\beta_\epsilon)\,\dd t$ and the exact power
in \eqref{rw:secondary:long-map}.
Appendix~\ref{rw:app:secondary-long} differentiates its renormalized
integral equation, including the moving endpoint and the height
$w_e=0$, and proves \eqref{rw:secondary:long-jets}.
The limiting unit retains its terminal-height dependence:
$\Psi_0(0,0,w_e;\theta)=-(2b)^{-1}\log(1+2w_e/k)$.
\end{proof}

The exact power can also be checked on the reversible field. Its
first integral is
$Z^{-2\theta}(X^2+\bar B^{-1}+2Z/k)$.
At fixed $Z/|X|$ its comparison with the maximum-height section gives
the power $b/k$, inverse to $\beta_\epsilon$.
For positive scale $1<b/k<2$, so that single power need not have
an ordinary second derivative at zero. This is why the long estimate
uses Euler input derivatives. Keeping $\beta_\epsilon$ exact is also
essential for unrestricted separation of scales: replacing it by its
zero-scale value would leave an error proportional to
$\epsilon\log H$, which need not be small.

\subsection{Two actual exits and compact section labels}

Locate the first exit of each returning upper arc to identify the middle
passage and define tags at its actual section points.

Fix a small $w_0>0$. At the wedge exit $w=w_0$, write
$\varphi=H$, $\gamma=\bar d/H$, $J=\bar m/H^2$.
The vertical passage \eqref{rw:secondary:vertical}, followed by
the reflected long map, has data
\[
 \widetilde H=HL_v,\qquad
 \widetilde\gamma=-\gamma/L_v,\qquad
 \widetilde J=J/L_v^2.
\]
Consequently the two original horizontal power inputs are
\begin{equation}\label{rw:secondary:wedge-inputs}
 u_\pm=H^{\beta_\epsilon}L_\pm,\qquad
 L_\pm,L_\pm^{-1}>c>0.
\end{equation}
Their logarithms have the finite jets above.
On a fixed original coefficient fibre the derivative is
\[
 \mathcal D=H\partial_H-\gamma\partial_\gamma-2J\partial_J,\qquad
 \mathcal D\log u_\pm
      =\beta_\epsilon+O(\delta+H^q)>c>0
\]
when $|\gamma|,|J|<\delta$.
Reflection gives $u_+=u_-$ on $\bar a=\bar d=\nu=0$.

For $\lambda>0$ define
\begin{equation}\label{rw:secondary:normalization}
 \lambda=\max(|\bar d|,\sqrt{|\bar m|}),\quad
 C=\bar d/\lambda,\quad M=\bar m/\lambda^2,\quad
 \max(|C|,\sqrt{|M|})=1,\quad
 \Lambda=\lambda^{\beta_\epsilon}.
\end{equation}
Coefficient differentiation is performed on ordinary open neighborhoods
of these signed faces, keeping $\lambda$ fixed rather than
redifferentiating the maximum.
Put $\hat x=\lambda X$, $\hat z=\lambda^2Z$,
$\delta_s=\epsilon\lambda^2$ and $\dd t_s/\dd\xi=1/\lambda$.
The exact field is
\begin{equation}\label{rw:secondary:field}
\begin{aligned}
 \dot{\hat x}&=(\hat z-\delta_s)[M(\hat z-\delta_s)-1]
                         +\theta\hat x^2+\epsilon\bar a\lambda\hat x,\\
 \dot{\hat z}&=\hat z[\hat x+C(\hat z-\delta_s)]
                         -\epsilon\nu\lambda(\hat z-\delta_s).
\end{aligned}
\end{equation}
There is no coefficient $\nu/\lambda$.
At zero scale division by $\hat z>0$ gives the affine field
\begin{equation}\label{rw:secondary:linear}
 \hat x'=M\hat z-1,\qquad \hat z'=\hat x+C\hat z.
\end{equation}

\begin{proposition}[First exits of returning upper arcs]
\label{rw:secondary:routing}
Consider actual upper arcs returning from $(p,0)$ to $(-q,0)$
with primary data tending to the stated corner, $\nu\to0$ and
$p\to\infty$. Closure at the central section is not required.
After fixed choices of sections and one common coefficient restriction,
each arc has either a first wedge exit
\[
 w=w_0,\quad 0<H<H_0,\quad
 |\bar d|<\delta H,\quad |\bar m|<\delta H^2,
\]
or a passage between $\hat\varphi=\hat x/\hat z=L$ and $-L$
in one fixed compact projective rectangle
$0<R=1/\hat z<R_{\max}$.
For $\lambda=0$ only the first alternative occurs.
No uniform flight-time bound is asserted.
\end{proposition}

\begin{proof}
The arrival estimate in Appendix~\ref{rw:app:secondary-barriers}
sends $p\to\infty$ to the fixed section $\varphi=\varphi_0$ with
$w\to0$, without any bound on $\epsilon\log p$.
Also $q\to\infty$: otherwise a bounded positive subsequence of $q$
and the reflected regular upper transition bound $p$; if $q\to0$,
the small-entry equation gives $p/q\to1$. Both contradict
$p\to\infty$. These arguments involve returning arcs, not their
central closing values.

Choose the wedge constants first, then $L\ge10$ so large that
$1/L<\delta/2$, $1/L^2<\delta/2$.
Keep the adjoining signs $W<-c$, $G+2W<-c$ on
$\lambda L\le\varphi\le\varphi_0$, $0<w\le w_0$.
Here $\varphi$ decreases and $w$ increases until a first terminal
side is reached. At $w=w_0$ the required ratios are at most
$1/L,1/L^2$. Its value $H$ tends to zero, since on every fixed
$\varphi\ge\varphi_1>0$ strip the time and multiplicative height
growth are bounded while the entry height tends to zero.
The other side is $\varphi=\lambda L$; assign a corner hit to
either alternative. If $\lambda=0$, exponential height growth
reaches $w=w_0$.

At the second side $0<R\le L^2w_0$.
The fixed first-exit barrier of
Appendix~\ref{rw:app:secondary-barriers} confines the arc until
$\hat\varphi=-L$. This crossing must occur for a returning arc,
whose incoming endpoint has
$\hat\varphi=-q/(\epsilon\lambda)<-L$.
It then has $w<R_{\max}/L^2<w_0+12$; reflected adjoining
monotonicity continues the same arc to the original horizontal
section. Uniqueness excludes another middle excursion.
The spatial barrier is not a claim that arbitrary incident
trajectories return.
\end{proof}

\begin{lemma}[Actual secondary tags]\label{rw:secondary:tags}
On the fixed outgoing secondary section, its actual point $(L,R)$
has a positive affine-flow tag $P\in[p_0,\infty]$.
The reciprocal tag $t=1/P$ is a regular increasing coordinate in $R$,
including $R=0$, with bounded finite jets and inverse jets.
The incoming section has the reflected tag $Q$ with the same
properties. These tags are not defined as $\lambda p,\lambda q$.
\end{lemma}

\begin{proof}
Follow \eqref{rw:secondary:linear} backward from the actual section
point. Multiplying by its fixed initial radial value $R$ gives
\[
 (\bar x,\bar z)(0)=(L,1),\qquad
 \dot{\bar x}=R-M\bar z,\quad \dot{\bar z}=-\bar x-C\bar z.
\]
Before $\bar z=0$, $0\le\bar z\le1$ and
$\bar x(s)\ge L-s$. A first-exit comparison gives a transverse
zero at $0<T<2/L$, with $\bar x(T)>L/2$.
This solution is smooth on the entire compact set
$0\le R\le R_{\max}$, $|C|,|M|\le1$.
Thus $P=\bar x(T)/R\ge L/(2R_{\max})=:p_0$ and
$t=R/\bar x(T)$ vanishes simply at zero.
Away from zero the flow determinant gives
\[
 \det(f,a_R)=\frac{R-M+L(L+C)}{R^3}>0,\qquad
 \det(f,(1,0))=-P<0
\]
at the two sections, respectively; hence $P_R<0$ and $t_R>0$.
Compactness supplies a common inverse margin. Reflection proves
the incoming statement.
\end{proof}

The actual section points and the first-exit choice give the compact data
\[
 (\bar B,\bar a,\epsilon,\lambda,C,M,t_P,t_Q),
 \quad \max(|C|,\sqrt{|M|})=1,\quad
 0\le t_P,t_Q\le p_0^{-1},
\]
together with the wedge or secondary exit choice and its fixed
section. At $\lambda=0$ the secondary entries are limiting
directions, not values of a divided original coefficient.
The same construction applies to nonperiodic returning arcs and
will be used for joint limits of an outer orbit and intermediate arcs.

\subsection{Exact cancellation and regular comparisons}

The original lower maps have one common horizontal exponent. Its
absolute power supplies the bound on $\nu$ at a closing zero;
cancellation in the derivative quotient supplies the relative expression
used to count zeros on the fixed fibre.

The original horizontal lower preparation is
\begin{equation}\label{rw:secondary:lower}
 D_\pm(u)=K_\pm u^\mu e^{P_\pm(u)/\epsilon},\quad
 P_\pm(0)=0,\quad \|P_\pm\|_{C^N}\le C,\quad
 H_\pm(u):=h_I+uP_\pm'(u)>c .
\end{equation}
This is \eqref{rw:transitions:horizontal-power}, in the two reflected
coordinates. A finite lower double root is crossed between fixed
slow sections before the final fixed core. Its whole-path condition
must hold; otherwise the passage is absent. On its passing side the
additional action belongs to $\epsilon\log K_\pm$ alone.
The relative normal estimate supplies the same ordinary input bounds
for $P_\pm$. No bounded coefficient derivatives of that divergent
action are inferred.

At $\hat\varphi=\pm L$ apply the variable-height long map, where
$H=\lambda L$, $\gamma=C/L$, $J=M/L^2$.
After regular outer connectors it gives
\begin{equation}\label{rw:secondary:normalized-inputs}
 u_\pm=\Lambda v_\pm .
\end{equation}
At a finite regular tag $v_\pm,v_\pm'>c$ after choosing a common
oriented input. At strict-focus infinity, using $t=1/P$,
$v_\pm=tV_\pm(t)$ with positive bounded units, derivatives and
inverse margins. The angular field there has a negative margin
on the compact passage, so its time and its ordinary radial jets
through $R=0$ are bounded. At singular tags the two outer maps
are constructed separately.

The long unit converges as $\lambda\to0$ uniformly in $\theta$;
then $\theta\to0$ identifies its limit with the affine tag flow:
\begin{equation}\label{rw:secondary:phase-ratio}
 v_-/v_+\longrightarrow P/Q(P;M,C)
\end{equation}
with the selected section and coefficient jets, including
strict-focus reciprocal zero. With $s$ denoting the long variable,
the affine backward leg has $\hat\varphi=Le^s$,
$\hat x=[Le^sw(s)]^{-1}$ and
$w(s)\sim w_e e^{-s}e^{\Psi_0}$. Its axis endpoint therefore satisfies
\[
 P^{-1}=Lw_e e^{\Psi_0(C/L,M/L^2,w_e;0)}.
\]
The original horizontal phase is $\varphi_0w_{\rm in}$ times a unit
tending to one, so its quotient by $\Lambda$ also tends to $P^{-1}$.
Reflection gives $Q^{-1}$ on the other leg.
The ordinary height estimates and the radial passage give the stated
jets after removing the common simple factor at reciprocal zero.

Let $I_\pm=D_\pm(u_\pm)$ and
$F=I_+-I_--\sqrt\epsilon\,\nu H_c$, where $H_c>c$ is independent
of input. For a common increasing input let
$l_\pm=D\log u_\pm$; use $D=\mathcal D$ in the wedge and
$D=t\partial_t$ on a secondary reciprocal interval.
The complete derivative logarithm is
\begin{equation}\label{rw:secondary:complete-log}
\begin{aligned}
 \mathcal K={}&\epsilon\log(I_-'/I_+')
 =\epsilon\log(K_-/K_+)+h_I\log(u_-/u_+)\\
 &+P_-(u_-)-P_+(u_+)+\epsilon\log(l_-/l_+)\\
 &+\epsilon\log\frac{H_-(u_-)}{H_+(u_+)} .
\end{aligned}
\end{equation}
In this quotient, cancel the identical power $\Lambda$, or
$H^{\beta_\epsilon}$, before any coefficient derivative; at reciprocal
zero cancel the common simple factor $t$ as well. The exact cancellation
removes $(\partial\beta_\epsilon)\log\lambda$ from the relative terms.
Single-sided compositions still contain factors
$\Lambda|\log\lambda|^j$, which a reserved positive power absorbs.
Thus the estimates allow unbounded $\epsilon|\log\lambda|$.

\begin{lemma}[Absolute scale budget]\label{rw:secondary:flat-budget}
On each regular-lower center comparison box, for arbitrary fixed
$A,N>0$, sufficiently small radial and scale cutoffs give
\[
 |I_+|+|I_-|\le Ce^{-A/\epsilon}\lambda^N
 \quad\hbox{or}\quad Ce^{-A/\epsilon}H^N
\]
with the finite coefficient derivatives used for center division.
The same absolute upper bound holds on a passing lower-double
overlay, whose extra action only decreases the amplitudes.
The derivatives of that extra action are excluded.
\end{lemma}

\begin{proof}
Before any scale factor is discarded, the exact combined power is
$\Lambda^\mu=\lambda^{k/(\epsilon\bar B)}$.
Bound the remaining amplitudes by $e^{C_0/\epsilon}$.
Choose $\lambda_0$ so that
$\frac12\inf(k/\bar B)\log(1/\lambda_0)>C_0+A$,
then choose $\epsilon$ so that
$\epsilon N<\frac12\inf(k/\bar B)$.
This proves the estimate. Further reserved powers absorb every
fixed logarithm, inverse scale loss and coefficient derivative.
The proof for $H$ is identical; at focus infinity reserve part
of the common input power as well. A fixed lower core has an
additional positive contraction budget, uniformly on compact
coefficient blocks, by Proposition~\ref{rw:relative-core}.
A fixed lower bottleneck contributes only further positive action.
\end{proof}

On a common regular-lower box reflection gives the exact identity
$I_+=I_-$ when $\bar a=\bar d=\nu=0$ in the wedge, or
$\bar a=C=\nu=0$ at secondary scale.
On that common box, divide the absolute difference by the transverse
coefficients using this identity. Only then impose the closing equation.
The absolute budget and Lemma~\ref{rw:centers:relative-forcing} give,
with $\eta_\epsilon\to0$,
\begin{equation}\label{rw:secondary:local-forcing}
 |\nu|\le\eta_\epsilon(|\bar a|+|\bar d|)
 \quad\hbox{in the wedge},\qquad
 |\nu|\le\eta_\epsilon\lambda^N(|\bar a|+|C|)
 \quad\hbox{at secondary scale}.
\end{equation}
The exponent $N$ is any preselected finite exponent after reserving
the required extra losses. In the first formula choose at least two
powers of $H$ before replacing $H^N|\bar d|/H$ by $|\bar d|$.
These scale powers come from the absolute deviations. Once a closing
zero has supplied the coefficient inequality, the critical-function
estimates use only the weaker bounds
\begin{equation}\label{rw:secondary:critical-forcing}
 |\nu|\le\eta_\epsilon(|\bar a|+|\bar d|)
 \quad\hbox{in the wedge},\qquad
 |\nu|\le\eta_\epsilon(|\bar a|+|C|)
 \quad\hbox{on regular secondary intervals},
\end{equation}
imposed on a fixed coefficient fibre, with $\eta_\epsilon\to0$.

For regular lower coefficients the endpoint constant is
\begin{equation}\label{rw:secondary:lower-constant}
 L(\bar a,\bar B)=
 -\frac{\pi\bar a}{\bar B\sqrt{4\bar B-\bar a^2}} .
\end{equation}
Indeed its limiting difference is
$-2\bar a\int_0^\infty X^2/(F_+F_-)\,\dd X$.
With $s=\sqrt{\bar B}X$, reciprocal symmetry of the quartic
reduces the integral to
$\pi/(2\bar B\sqrt{4\bar B-\bar a^2})$.
The subtracted tail is integrable with coefficient derivatives on
every compact subset of $|\bar a|<2\sqrt{\bar B}$.
The exact horizontal power subtraction and the finite-order lower
phase estimates transfer this limit with the selected jets.
At a lower double root, instead, retain
$\mathcal K=\mp A_l+O(1)$ on its whole-path passing set;
the sign corresponds to the contracting side. It is a strict
constant-sign comparison, not a bound on $\mathcal K'$.

\begin{proposition}[Wedge and regular secondary rows]
\label{rw:secondary:regular}
\label{rw:secondary:wedge}
The wedge has local bound $Z(F)\le2$, and so does every finite
regular secondary tag or strict-focus infinite tag.
Each assertion holds on an open coefficient and section neighborhood,
includes both signs of $\nu$, and retains all lower passing conditions.
Near a regular center the derivative logarithm has at most one zero
under the respective bound in \eqref{rw:secondary:critical-forcing}.
For the closing-zero
claim that bound follows from any one closing zero; if none exists,
there is nothing to count.
A lower-double action-dominated neighborhood has a strict sign of
$\mathcal K$ and hence at most one closing zero.
\end{proposition}

\begin{proof}
In the wedge, divide \eqref{rw:secondary:complete-log} on its
common center box:
\begin{equation}\label{rw:secondary:wedge-columns}
 \mathcal K=\bar a A+(\bar d/H)C_w+\nu E,\qquad
 A\to-\frac{\pi}{2\bar B^{3/2}},\quad
 C_w\to-\frac{2c_w}{3\bar B},\quad
 c_w=\frac{\sqrt{1+2w_0}}{w_0}.
\end{equation}
The $\mathcal D$ derivatives of these columns tend to zero.
To compute the second column, the limiting center parabola is
$Z=(p^2-X^2)/2$ and $p=c_w/H$ at $w=w_0$.
Variation of $Z''-\bar dZ'=-1$ gives
$\partial_{\bar d}(q-p)=2p^2/3$.
Its endpoint action contribution is
$-2Hp^3/[3(1+\bar Bp^2)]$, with the displayed limit.
The first column is the derivative of
\eqref{rw:secondary:lower-constant} at zero.
The fixed-section endpoint variation cancels between the two
center maps. The long estimates transfer these calculations with
their $\mathcal D$ jets on small $|\gamma|,|J|,H,\epsilon$.

Absorb the forcing column by adding
$\nu\bar a E/(\bar a^2+\bar d^2)$ to $A$ and
$H\nu\bar d E/(\bar a^2+\bar d^2)$ to $C_w$.
The fixed-fibre operator has
$\mathcal D\bar a=\mathcal D\bar d=\mathcal D\nu=0$.
Thus \eqref{rw:secondary:critical-forcing} makes the corrections
$O(\eta_\epsilon)$ and $O(H\eta_\epsilon)$ in the required
$\mathcal D$ jets. Both columns retain negative
margins and
$\mathcal D\log[C_w/(HA)]<-1/2$.
Thus their ratio is strictly monotone and $\mathcal K$ has at most
one zero. Zero transverse coefficients give the exact center.
Strict nonzero $\bar a$ and the lower-double boundary are instead
covered by the value or action sign in
\eqref{rw:secondary:lower-constant}. Compactness of the closed
lower cone supplies a finite choice of such neighborhoods; one
neighborhood counts the entire smaller wedge of a fixed fibre.

At regular secondary tags the limit is
\begin{equation}\label{rw:secondary:regular-limit}
 \mathcal K_0=L(\bar a,\bar B)-\bar B^{-1}\log(Q/P).
\end{equation}
We record the affine-flow facts that determine its columns.
For return time $T$, set $\zeta(s)=e^{-Cs/2}\hat z(s)$ along the
orbit.
If $S_\chi''=\chi^2S_\chi$, $S_\chi(0)=0$, $S_\chi'(0)=1$,
$\chi^2=M+C^2/4$, the Dirichlet formula gives
\[
 P=\frac{\int_0^TS_\chi(T-s)e^{-Cs/2}\,\dd s}{S_\chi(T)},\quad
 Qe^{-CT/2}=\frac{\int_0^TS_\chi(s)e^{-Cs/2}\,\dd s}{S_\chi(T)}.
\]
The first-return time precedes the first positive zero of $S_\chi$.
Pair $s$ and $T-s$ for $0<s<T/2$: the kernel difference
$S_\chi(T-s)-S_\chi(s)$ is positive by the hyperbolic sine formula
or its linear limit, and for $\chi=i\omega$ it equals
$2\cos(\omega T/2)\sin(\omega(T/2-s))/\omega>0$
because the first-return condition gives $T<\pi/\omega$.
On this half-interval, $e^{-Cs/2}-e^{-C(T-s)/2}$ has sign $C$,
so the paired integral gives
$\operatorname{sign}(P-Qe^{-CT/2})=\operatorname{sign}C$.
Together with $Q_P=(P/Q)e^{CT}$ this proves
$\operatorname{sign}\partial_P\log(Q/P)=\operatorname{sign}C$.
On $C=0$ the divided column is
\[
\begin{gathered}
 G(P,M)=\partial_C\log(Q/P)|_{C=0}
   =\frac2{P^2}\int_0^{z_{\max}}
                         \sqrt{P^2-2z+Mz^2}\,\dd z,\\
 G_P=2\int_0^1\frac{v^2}{1-MP^2v^2}\,\dd v>0.
\end{gathered}
\]
The vanishing endpoint integrand justifies its differentiation.
For $M=-r^2$, direct integration gives
$G(1/t,M)=\pi/(2r)-2t/r^2+O(t^2)$.
More generally at strict-focus infinity the affine solution gives
\[
 \log(Q/P)|_{t=0}=\frac{\pi C}{\sqrt{-4M-C^2}},\quad
 \partial_t\log(Q/P)|_{t=0}
 =-\frac{C(1+e^{-\pi C/\sqrt{-4M-C^2}})}{-M}.
\]
These follow by imposing the first zero in
\[
 \hat z(s)=\frac1M\left[
 1-e^{Cs/2}\left(\cos\omega s-
       \frac{MP+C/2}{\omega}\sin\omega s\right)\right],
 \qquad \omega=\sqrt{-M-C^2/4};
\]
they have finite coefficient jets on strict-focus compact sets.

For $C\ne0$ these derivative margins transfer to $\mathcal K$.
For $C=0,\bar a\ne0$ use its nonzero value.
Near $\bar a=C=0$, divide into
$\mathcal K=\bar a A+CD+\nu E$; the limiting columns are
$A_0=-\pi/(2\bar B^{3/2})$ and
$D_0=-G(1/t,M)/\bar B$.
Their ratio has a strict derivative margin, also at $t=0$.
The mixed jets of Proposition~\ref{rw:exact-long} and the horizontal
preparation give $\|E\|_{C^1}\le C$ uniformly on this regular-lower
box. To see this, cancel $\Lambda$ in
\eqref{rw:secondary:complete-log} before taking one $\nu$ and one
input derivative. At focus infinity write $v_\pm=tV_\pm$ and cancel
$t$ as well. The remaining relative units have positive margins;
one input derivative of their logarithmic rates and one coefficient
derivative spend two input orders and one coefficient order, with
one further input order reserved for the simple-factor division.
At fixed $\lambda,C$, the original lower coefficient is
$\bar c=\lambda C-\nu$, so the $\nu$ derivative is
$\partial_\nu-\partial_{\bar c}$ in the original lower box.
The single-sided terms retain a positive scale power.
For fixed horizontal and central sections, shrink the regular center box
so that $|\bar a|\le\sqrt{B_*}/2$ and $\epsilon|\bar m|\le1/4$.
The exact formula \eqref{rw:transitions:lower-factor}, with a smaller
scale cutoff, gives $F_I^\pm(s)\ge B_*/2+s^2/2$ on $[0,x_R]$.
Write $\beta_I^\pm/(sF_I^\pm)=h_I/s+q_\pm$ and use the fixed
split $x_h$ of Proposition~\ref{rw:lower-preparation}. The horizontal
proof and the curved and central connections give
\[
\begin{aligned}
 \epsilon\log K_\pm={}&-S_c^\pm-h_I\log x_h
 -\int_0^{x_h}q_\pm(s)\,\dd s
 -\int_{x_h}^{x_R}\frac{\beta_I^\pm(s,0)}{sF_I^\pm(s)}\,\dd s\\
 &+\epsilon\log U_\pm-\epsilon\log\epsilon,
 \qquad U_\pm=Z_e^\pm(0)e^{j_\pm(0)}u_c^\pm(0)a_c^\pm.
\end{aligned}
\]
The units $U_\pm$ are positive with bounded logarithmic coefficient
jets by Propositions~\ref{rw:transitions:curved-exit}
and~\ref{rw:transitions:central-core}. For
$\mathscr D_\nu=\partial_\nu-\partial_{\bar c}$, fixed $R$ gives
$\mathscr D_\nu h_I=\mathscr D_\nu S_c^\pm=0$.
Differentiating the fixed integrals using the displayed positive
denominator and Proposition~\ref{rw:phase} therefore bounds
$\mathscr D_\nu[\epsilon\log(K_-/K_+)]$ uniformly on this regular
center box. Define the forcing column by integrating this
$\nu$ derivative from zero. The corrections
\[
 \frac{\nu\bar a E}{\bar a^2+C^2},\qquad
 \frac{\nu C E}{\bar a^2+C^2}
\]
are therefore $O(\eta_\epsilon)$ in $C^1$ under
\eqref{rw:secondary:critical-forcing}. Adding them to $A,D$
preserves their margins and the strict derivative of their quotient.
This gives $Z(\mathcal K)\le1$ and $Z(F)\le2$.
The lower-double case uses its independent constant action sign.
\end{proof}

\begin{lemma}[Outward direction and regular departure]
\label{rw:secondary:chain-direction}
\label{rw:secondary:chain-monotone}
The wedge's outward derivative is a positive multiple of
$-\mathcal D$, and the secondary tag $P$ increases outward.
Under the corresponding relative forcing bounds, if
$\bar a\bar d<0$ and $\mathcal K=0$ or
$\operatorname{sign}\mathcal K=-\operatorname{sign}\bar d$,
its outward derivative has sign $-\operatorname{sign}\bar d$
on the wedge and on a fixed compact regular secondary set.
\end{lemma}

\begin{proof}
For the actual input $a(p)=(p,\epsilon)$ and wedge output
$b(H)=(1/(w_0H),1/(w_0H^2))$, the flow determinants are
$-\epsilon p<0$ and $\dot w/(w_0^3H^4)>0$.
The flow Jacobian has positive determinant, so $H'(p)<0$ and
$\dd/\dd p=(H'(p)/H)\mathcal D$ exactly.
The tag orientation follows from Lemma~\ref{rw:secondary:tags}
and the ordering of the transverse outgoing intersections.

Put $s=\operatorname{sign}\bar d$.
The hypothesis $s\mathcal K\le0$ in the wedge first excludes the
strict nonzero lower-cone values and then, by its two negative
columns, bounds $|\bar a|\le C|\bar d|/H$.
Inserting this bound into $-s\mathcal D\mathcal K$ leaves the
strict negative main term $(|\bar d|/H)C_w$; the column derivatives
are arbitrarily smaller. At secondary scale the same hypothesis
and \eqref{rw:secondary:regular-limit} give $|\bar a|\le C|C|$.
The outward derivative of $\log(Q/P)$ has sign $C$ and a margin
relative to $|C|$, using $-t$ as the regular outward coordinate at
focus zero.
The divided relative errors are small compared with that margin.
Thus the outward derivative of $\mathcal K$ has sign $-s$.
This is a local direction statement on actual open regular overlaps;
it does not add their zero counts.
\end{proof}

\subsection{Singular secondary endpoints}

For use at every singular row keep the exact projective equation.
With $R=1/\hat z$, $\hat\varphi=\hat x/\hat z$,
$A_s=\epsilon\bar a\lambda$, $\nu_s=\epsilon\nu\lambda$, it is
\begin{equation}\label{rw:secondary:projective}
\begin{aligned}
 R'={}&-R(\hat\varphi+C)+(C\delta_s+\nu_s)R^2-\nu_s\delta_sR^3,\\
 \hat\varphi'={}&M-b\hat\varphi^2-C\hat\varphi
 +R[-1-2M\delta_s+(A_s+C\delta_s+\nu_s)\hat\varphi]\\
 &+R^2[\delta_s+M\delta_s^2-\nu_s\delta_s\hat\varphi].
\end{aligned}
\end{equation}
Its coefficients converge on fixed compact rectangles. The lower
coefficient $\bar B$, however, is still measured at the original scale;
replacing it by $\bar B/\lambda^2$ would lose the compact coefficient
range required by the lower estimates. We therefore attach each new
upper map to the original lower maps \eqref{rw:secondary:lower}.

\begin{proposition}[A finite secondary saddle]
\label{rw:secondary:saddle}
For $M>0$, the limiting tags
$P=P_c(M,C)$, $Q=P_c(M,-C)$ have local bound $Z(F)\le3$.
Let $\kappa=r-1$ for the actual saddle ratio $r$.
Under $|\nu|\le\eta_\epsilon(|\bar a|+|\kappa|)$, with
$\eta_\epsilon\to0$, their derivative logarithm also satisfies
$Z(\mathcal K)\le2$.
Here
\[
 P_c(M,C)=\frac2{C+\sqrt{C^2+4M}} .
\]
The actual domain is the positive returning saddle quadrant, intersected
with every lower whole-path condition and the positive central outputs.
A lower double root may shorten this domain.
\end{proposition}

\begin{proof}
The limiting saddle is $(-C/M,1/M)$, with ratio
$r_0=(\sqrt{C^2+4M}-C)/(\sqrt{C^2+4M}+C)$.
Near its positive height, division by $\hat z$ turns
\eqref{rw:secondary:field} into an analytic $O(\epsilon)$
perturbation of the affine field, uniformly through $\lambda=0$.
The axis construction of Section~\ref{rw:upper-resonance} therefore
applies on a common neighborhood. Its two outer maps are
$u_+(t)=\Lambda v_+(t)$, $u_-(y)=\Lambda v_-(y)$ with
$v_\pm,v_\pm'>c$; cancellation of $\Lambda$ gives
$v_-(0)/v_+(0)\to r_0$.

Near $C=0$ preserve the exact reversible identity.
On $\bar a=C=\nu=0$ put $z_s=\delta_s+1/M$ and
\[
 W_c(z)=(z/z_s)^{-2\theta},\qquad
 I_c(z)=\int_{z_s}^z
       \frac{2(s-\delta_s)[M(s-\delta_s)-1]}sW_c(s)\,\dd s .
\]
The first integral is $W_c\hat x^2-I_c$.
Since $I_c=(z-z_s)^2J_c$ with $J_c(z_s)=1/z_s>0$,
its analytic square root supplies reflected product coordinates.
Extend them and normalize the neighboring axes, using the identity
when they are already normalized. The reciprocal preparation then
gives the actual ratio $\kappa=r-1$ and
\[
 y=t^re^w,\qquad
 w=-\sigma_s t\omega(t;\kappa)+O_E(\epsilon n t^{7/4}),
 \qquad |\sigma_s|\le C\epsilon n,\quad
 n=|\bar a|+|\kappa|+|\nu|.
\]
The identity is $y=t$ on the center. Hadamard division on the
common coefficient box supplies the factor $n$ in the remainder.
Finite mixed jets are obtained from the same equation and
coefficient division on its fixed positive quadrant.

Put $U=D_+(u_+)$, $V=D_-(u_-)$ and
$\mathsf A_\pm=\epsilon U''/U',\epsilon V''/V'$.
For $u=\Lambda v$, set
$Z_\pm=H_\pm(u)v'/v$.
Then $\mathsf A_\pm=Z_\pm+\epsilon Z_\pm'/Z_\pm$.
These rates have positive margins and bounded ordinary jets.
A lower-double action is a constant multiplier and cancels here.
It is not differentiated in coefficients.

Integrating these rates gives \eqref{rw:finite:saddle-log} with
$(v,K)=(y,\mathcal K)$:
\begin{equation}\label{rw:secondary:saddle-log}
 \mathcal K=K_c+\epsilon\kappa\log t+\mathcal E(t),\qquad
 K_c=\epsilon\log[rV'(0)/U'(0)] .
\end{equation}
The remainder and its Euler derivative are $O(t^{1/2})$,
or $O(nt^{1/2})$ near the center. Cancellation of $\Lambda$
before differentiating shows
\[
 K_c\longrightarrow L(\bar a,\bar B)+\bar B^{-1}\log r_0 .
\]
On a regular center box this yields
$K_c=\bar a A_c+\kappa B_c+\nu C_c$, where
$A_c\to-\pi/(2\bar B^{3/2})$ and the other columns are bounded.
The exact center and the absolute budget give
$|\nu|\le\eta(|\bar a|+|\kappa|)$ at a closing zero.
The actual ratio coordinate is legitimate because
$\partial_C\kappa\to-1/\sqrt M$.

We apply Lemma~\ref{rw:finite:saddle-counting} to this preparation.
For a nonzero limiting $\bar a$, clause \textup{(i)} uses $R_*=1$;
for a dominating lower action it uses that action's magnitude.
Near the center fix a large $K_0$. The divided constant and relative
forcing give the same hypotheses with $R_*=|\bar a|$ in the cone
$|\bar a|>K_0|\kappa|$.

In the remaining cone the constant may cancel, so we count the Euler
derivative instead. For $|\bar a|\le K_0|\kappa|$, the coefficient
bound gives $n\le(1+K_0)(1+\eta_\epsilon)|\kappa|$.
If $\kappa=0$, it first gives $\bar a=\nu=0$, the exact center.
On the noncenter part, differentiation gives
\eqref{rw:finite:saddle-compensator} with $K=\mathcal K$,
$A=\mathsf A_-(0)$, $B_0=\mathsf A_+(0)$ and $\sigma=\sigma_s$;
the prepared remainder satisfies $|D^jR|\le Cnt^{3/2}$, $j\le2$.
Since $\sigma_s/\kappa=O(\epsilon)$, the positive rate margin gives
$W/\kappa>c$, $|T|\le C|\kappa|$ and
$|R'|=|DR|/t\le C|\kappa|t^{1/2}$.
Clause \textup{(ii)} gives $Z(\mathcal K)\le2$.
Away from ratio one, the unsheared reciprocal equation gives
$y=k_0t^re^w$ with a bounded positive $k_0$ and the derivative
expansion of clause \textup{(iii)}, with
$A=\mathsf A_-(0)$ and $B=\mathsf A_+(0)$.
Thus the same bound holds there. By \ref{rw:zero-calculus},
$Z(F)\le3$ on the lower-passing, positive-central subinterval;
no endpoint value at a truncated passage is used.
\end{proof}

We also retain the actual coefficient identity needed by a later
entrance argument. On $\bar a=\nu=0$ let $q=M+\theta C^2$.
The saddle Jacobian of the divided field is exactly
\[
 J_s^{\rm ref}=
 \begin{pmatrix}
 -2\theta C/(1+\delta_s q)&(M-\theta C^2)/(1+\delta_s q)\\
 1&C
 \end{pmatrix}.
\]
The remaining coefficients change it by
$O(\epsilon\lambda(|\bar a|+|\nu|))$.
Thus on a compact positive-$M$ block
\begin{equation}\label{rw:secondary:saddle-ratio}
 \kappa=-\alpha_s C+
 O(C^2+\epsilon\lambda|\bar a|+\epsilon\lambda|\nu|),\qquad
 \alpha_s=\frac{1+M\delta_s-2\theta}
                  {\sqrt{M(1+M\delta_s)}}>c>0 .
\end{equation}
The reference ratio is exactly one at $C=0$. This is an identity
for the secondary parameters $C=\bar d/\lambda$; its small
$\nu$ term does not justify dropping the order-one $\nu$ term in
the primary ratio when $\bar c$ is held fixed.

\begin{proposition}[Real, repeated and weak secondary rows]
\label{rw:secondary:real}
For $C>0$ the following limiting rows have the stated local bounds;
reflection supplies $C<0$:
\[
\begin{array}{c|c|c}
\text{row}&\text{limiting tags}&Z(F)\\ \hline
 M<0,\ C^2+4M>0&P=P_c,\ Q=\infty&\le2\\
 C=2g,\ M=-g^2&P_c=1/g\le P<\infty,\ Q=\infty&\le2\\
 C=2g,\ M=-g^2&P=Q=\infty&\le2\\
 M=0&P=1/C,\ Q=\infty&\le3
\end{array}
\]
Every assertion holds in an open coefficient and section neighborhood.
It includes neighboring signs of the unfolding parameters and either
lower double root on its whole-path passing set. Positive central
outputs are imposed in addition to the upper conditions below.
\end{proposition}

\begin{proof}
In every row the lower maps remain \eqref{rw:secondary:lower}, with
divergent lower actions confined to input-independent constants.
The upper singularities require different estimates: a hyperbolic power
for the strict row, a quadratic angular passage for the repeated rows,
and the exact moving radial factor for the weak row. We keep the passing
domains throughout these preparations and their adjacent unfoldings.

For the strict row define
\[
 \Delta=\sqrt{C^2+4bM},\quad
 \varphi_s=(-C-\Delta)/(2b),\quad
 \gamma=\varphi_s+C=\frac{Ck-\Delta}{2b},\quad r=\gamma/\Delta.
\]
Both $\Delta,\gamma$ have positive margins.
The projective saddle has eigenvalues $-\gamma,\Delta$.
Straighten its axes and integrate its orbital reciprocal equation
in the returning quadrant. The map is $k_0t^re^{j(t)}$,
with three selected Euler derivatives of $j$ tending to zero
as the section shrinks. Compact positive-ratio margins absorb
the fixed logarithmic losses, including internal resonances.
The other quadrant enters the attracting node.
After the two long legs,
\[
 u_+=\Lambda v(t),\quad v,v'>c,\qquad
 u_-=\Lambda k_1t^re^{j_1(t)} .
\]
In \eqref{rw:secondary:complete-log},
$t\partial_t\log l_+=1+O(t)$ and
$t\partial_t\log l_-=o(1)$.
Consequently $t\mathcal K'=h_Ir-\epsilon+o(1)>0$.
The terms from $P_\pm$ are $O(\Lambda t)+O(\Lambda t^{r_*})$.
This proves the two-zero closing bound.

Near a repeated root the boundary restriction of
\eqref{rw:secondary:projective} gives exactly, after the positive
division by $b$ and a finite-order phase projection,
\begin{equation}\label{rw:secondary:repeated}
\begin{gathered}
 v'=\mu_u+v^2=:f(v),\qquad R'=-R\beta_u(v,R),\\
 \beta_u(v,0)=(\Xi-v)/b,\qquad
 \mu_u=-\frac{C^2+4bM}{4b^2},\qquad \Xi=\frac{Ck}{2b}.
\end{gathered}
\end{equation}
The normal margin is positive. Fix the finite derivative order
before taking the angular strip small enough that its slow variation
growth is dominated by that margin, as in
Section~\ref{rw:upper-degenerate-preparation}.
The phase construction preserves this exact boundary equation.
A moving entry $v$ can reach a fixed positive exit $v_R$ precisely when
\begin{equation}\label{rw:secondary:repeated-domain}
 f(s)>0\quad(v\le s\le v_R):
 \quad \mu_u>0\ \text{or}\ %
       [\mu_u\le0,\ v>\sqrt{-\mu_u}].
\end{equation}
No scalar root can be crossed by using a different normal fibre.

For a finite repeated tag the entry coordinate decreases strictly
with $P$. At the limiting field this follows from the coordinates
$V,Z_1$ and $r'=-r^2$ in
Appendix~\ref{rw:app:secondary-tails}; the fixed radial entry
and transverse outer map preserve a negative derivative margin.
The normal passage has
\[
 G(v)=k_0\exp[-\int_v^{v_R}\beta_u(s,0)/f(s)\,\dd s]e^{j(v)},
 \qquad N=fG'/G=\beta_u(v,0)+fj'>c .
\]
The selected ordinary jets of $j$ are bounded by
Proposition~\ref{rw:relative-core} and
Lemma~\ref{rw:transitions:moving-rate}.
Write $u_+=\Lambda w(v)$, $w,w'>c$,
$u_-=\Lambda k_1G(v)$ and $\psi=w'/w$.
The full numerator is
\begin{equation}\label{rw:secondary:finite-repeated-numerator}
\begin{aligned}
 f\mathcal K'={}&h_IN-h_If\psi
       +u_-P_-'N-fu_+P_+'\psi\\
 &+\epsilon(fN'/N-f'-f\psi'/\psi)
       +\epsilon f\partial_v\log[H_-(u_-)/H_+(u_+)] .
\end{aligned}
\end{equation}
The last term is $O(\epsilon\Lambda)$; the other perturbations
are $O(\Lambda+|f|+\epsilon|f'|+\epsilon|f|)$.
First reduce the angular strip, then the scale, so $h_IN$
dominates. Thus $f\mathcal K'>0$ on the whole actual branch.

For the two-infinite repeated row, fixed slow sections
$-v_R,v_R$ require $\mu_u>0$. Their contraction is
\[
 e^{-A_u},\qquad
 A_u=\frac{2\Xi}{b\sqrt{\mu_u}}\arctan(v_R/\sqrt{\mu_u}).
\]
The outer relative units have ordinary radial jets, giving
$u_+=\Lambda tV_+(t)$,
$u_-=\Lambda e^{-A_u}tV_-(t)$ with $t=1/P$ and positive units.
Hence
\[
 \mathcal K(0)=\epsilon\log(K_-/K_+)-h_IA_u+O(1),\qquad
 \|\mathcal K-\mathcal K(0)\|_{C^2}\le C .
\]
The explicit repeated first integral and the uniformly contracting
tails of Appendix~\ref{rw:app:secondary-tails} prove
\begin{equation}\label{rw:secondary:repeated-limit}
 \mathcal K_P\longrightarrow
 \bar B^{-1}[P^{-1}-P/(P-P_c)^2],\qquad
 \mathcal K_t\longrightarrow
 \frac{P_c(2-P_ct)}{\bar B(1-P_ct)^2}>0 .
\end{equation}
The ordinary $C^2$ bound transfers this derivative limit to $t=0$:
compare first with a fixed small positive $t$ and then let that input
decrease to zero. The positive derivative gives the two-zero closing
bound even when upper and lower divergent actions cancel in the constant.

For the weak row take $C>0$; reflection is applied to the whole
field before choosing the contracting chart when $C<0$.
Put $\Delta=\sqrt{C^2+4bM}$ and
$\varphi_s=(-C-\Delta)/(2b)$. The exact radial eigenvalue is
\begin{equation}\label{rw:secondary:weak-face}
 \gamma=\varphi_s+C
 =-\frac{2(M+\epsilon\bar BC^2)}{Ck+\Delta}.
\end{equation}
Thus the moving weak face is $M+\epsilon\bar BC^2=0$.
In reversed projective time the normal eigenvalue is $-\Delta$.
The weak graph gives a reduced equation $\zeta'=\zeta G(\zeta)$,
where $G(0)=\gamma$ and
\[
 G'(0)=
 \frac{1+2M\delta_s-(A_s+C\delta_s+\nu_s)\varphi_s}
      {\Delta+\gamma}-C\delta_s-\nu_s\longrightarrow1/C>0.
\]
Choose the finite smooth complete extension and phase projection as
in Section~\ref{rw:upper-degenerate-preparation}. Its inputs hold here:
$R=0$ is invariant, the normal margin is $\Delta$, and the displayed
positive $G'(0)$ gives
$u=G(\zeta)-\gamma=\zeta\int_0^1G'(s\zeta)\,\dd s$.
Division by its positive time factor gives exactly
$u'=f(u)=u(\gamma+u)$, spending one radial derivative.
The derivative orders precede the weak width, so normal decay
dominates every retained slow variation.
The unique positive component reaching the fixed exit $u_R$ is
\begin{equation}\label{rw:secondary:weak-domain}
 \max(0,-\gamma)<u<u_0<u_R/2,\qquad |\gamma|<u_R/2.
\end{equation}
Intersect this domain with the two original lower whole-path
conditions and positive central outputs.

Let $t=t(u)$ be the inverse normal passage with its regular connector.
Proposition~\ref{rw:relative-core} and
Lemma~\ref{rw:transitions:moving-rate} give
\[
 \frac{t'}t=\frac Nf,\qquad
 \frac{t''}{t'}=\frac{A_U}f,\qquad
 A_U=N-f'+fN'/N,\qquad N,A_U>c,
\]
with bounded ordinary $C^2$ rates. Moreover $\|t\|_{C^3}\to0$
uniformly as $u_0,|\gamma|\to0$, with $u_R$ fixed:
the transit time tends uniformly to infinity, and
$f(u_R)/f(u)\le e^{\varpi T}$ when $|f'|\le\varpi$,
so the reserved normal exponential absorbs every fixed inverse-speed
loss. In the secondary upper construction retain joint order twelve
before the positive coordinate/time change, leaving eleven;
relative-core with $N=7$ uses nine primitive orders and gives seven
orders for the relative logarithm, hence five for its rates, including
two ordinary input and three coefficient orders.
Independently retain original-lower field, reference and connection
order ten for the Gaussian divided-$\nu$ bound in
Appendix~\ref{rw:app:relative-core}. The lower-phase reserve gives
raw graph order seventeen for this target, whereas the full lower-preparation reserve in
Proposition~\ref{rw:lower-preparation} gives eighteen
for $P_\pm$ through order seven. Choose the original lower graph
at order eighteen. Retain order seven also in $\Phi,w$ for the
rate and simple-factor divisions below. All orders count ordinary
section inputs, coefficients and Euler scale derivatives together.

Write $u_+=\Lambda w(t)$, $w,w'>c$, and
$u_-=\Lambda\Phi(u)$, $\Phi(0)=0$, $\Phi'>c$.
Put $L_u=u\Phi'/\Phi$, $R_\Phi=\Phi''/\Phi'$,
$\psi=w'/w$, and
\[
\begin{aligned}
 A_I&=H_-(u_-)-\epsilon+
       \epsilon u_-H_-'(u_-)/H_-(u_-),\\
 R_{\rm out}&=H_+(u_+)\psi+
       \epsilon[\Lambda w'H_+'(u_+)/H_+(u_+)+\psi'/\psi].
\end{aligned}
\]
The complete derivative quotient, with its common exact scale power
already canceled, gives
\begin{equation}\label{rw:secondary:weak-numerator}
 \mathcal M=f\partial_u\mathcal K
 =A_IL_u(\gamma+u)+\epsilon fR_\Phi
       -R_{\rm out}(t)Nt-\epsilon A_U.
\end{equation}
Here $A_I(0)=h_I-\epsilon=k/\bar B$,
$A_I',A_I''=O(\Lambda)$ and $L_u=1+O(u)$.
Lemma~\ref{rw:weak:rate-calculus} applies with
$h=1$, $A=A_IL_u$, $R=R_\Phi$ and $E=R_{\rm out}N$.
Their bounded ordinary $C^1$ rates and $\|t\|_{C^2}\to0$ give
$\partial_u\mathcal M>c$ after the weak ranges and then the scales
are reduced, retaining the term $-\epsilon A_U'$.
Hence $Z(\mathcal K')\le1$, $Z(\mathcal K)\le2$ and
$Z(F)\le3$ by \ref{rw:zero-calculus}. Lower-double constants
do not enter these input rate derivatives.

Each numerator argument uses only input derivatives of lower
bottleneck maps, so their divergent constants are harmless.
Reflection exchanges both actual maps and the oriented inputs,
takes $(\bar a,C,\nu)$ to their negatives, and preserves return
multiplicity. It is applied before choosing the contracting chart.
\end{proof}

The weak numerator supplies a useful direction without taking a
formal endpoint value: in the outward orientation $u$ decreases,
so $\operatorname{sign}\mathcal K'_{\rm out}
=-\operatorname{sign}\mathcal M$.
This remains valid on a shortened lower-passing interval.
For finite saddles \eqref{rw:secondary:saddle-ratio} supplies the
actual ratio direction. These directions give the entrance counts in
Propositions~\ref{rw:secondary:saddle-anchor}
and~\ref{rw:secondary:weak-anchor}.

\subsection{Exhaustion of the normalized escape faces}

It remains to show that these rows cover every returning-arc limit.
We use the compact actual tags to include the repeated and weak
boundaries, without assuming that the arc closes at the central section.

\begin{proposition}\label{rw:secondary:exhaustion}
Every limit of nonwedge returning upper arcs in
Proposition~\ref{rw:secondary:routing}, with an inherited lower cone,
belongs to one of the following rows or its reflection:
\[
\begin{array}{c|c}
 \text{limiting upper data}&\text{tags}\\ \hline
 \text{regular}&0<P<P_c,\quad Q=Q(P;M,C)\\
 M>0& P=P_c,\quad Q=Q_c\\
 M<0,\ C>0,\ C^2+4M>0&P=P_c,\quad Q=\infty\\
 C>0,\ C^2+4M=0&P_c\le P<\infty,\quad Q=\infty\\
 M=0,\ C>0&P=1/C,\quad Q=\infty\\
 M<0,\ C^2+4M<0&P=Q=\infty\\
 M<0,\ C^2+4M=0&P=Q=\infty .
\end{array}
\]
Here $Q_c=P_c(M,-C)$ and $P_c=\infty$ when
$C^2+4M<0$ or $C+\sqrt{C^2+4M}\le0$.
The secondary signed normalization remains in every row.
No third occurrence of the same zero corner is possible.
The assertion applies to returning arcs without a closing condition.
\end{proposition}

\begin{proof}
Take a convergent subsequence in the compact normalized face,
inherited lower cone and reciprocal tags. If $P_*<P_c$, the affine first return
is regular and compact finite-time convergence gives
$Q_*=Q(P_*;M,C)$. Reflection gives the converse when $Q_*<Q_c$.

Let $D=C^2+4M>0$ and $P_c<\infty$. A supercritical
$P_*>P_c$, including infinity, enters the boundary node
\[
 R=0,\qquad \hat\varphi_n=(-C+\sqrt D)/2,
\]
whose eigenvalues are $-\sqrt D$ and $-(C+\sqrt D)/2<0$.
For a finite tag the affine solution has a positive growing
coefficient; for an infinite tag follow the boundary angular equation
from $L$. A small trapping neighborhood in $R\ge0$ persists by its
strict eigenvalue margins and invariant boundary. Continuous
dependence on the preceding compact segment sends nearby arcs into
it and excludes arrival at $-L$, contradicting their return.
Reflection gives the opposite-end exclusion. This requires $D>0$.

If both finite tags are nonregular, the affine eigenline classification
therefore gives the two finite thresholds with $M>0$.
If just one tag is infinite, it gives the strict-real, finite-input
repeated, or weak row. At a repeated root the full finite interval
$P\ge P_c$ remains. A finite opposite tag would be a regular
reflected return and would force $P<P_c$.
If both tags are infinite, the node exclusions force $D\le0$.
Normalization then implies $M<0$, since $M=0$ would force
$C=0$ and violate $\max(|C|,\sqrt{|M|})=1$.
These give precisely the two-infinite rows.

On the normalized repeated face $|C|=1$, $M=-1/4$;
on the limiting weak face $|C|=1$, $M=0$.
Their adjacent signed unfoldings have the whole-path domains
\eqref{rw:secondary:repeated-domain} and
\eqref{rw:secondary:weak-domain}; the reciprocal repeated tag has
its separate relative preparation. Together with the strict-real,
finite-saddle and focus rows, these cover every boundary of the
normalized face, with actual first-exit flags and tags retained.
Normalization excludes $(C,M)=(0,0)$, so another parabolic zero
corner cannot occur.

An excluded returning sequence at arbitrarily small scales would
contradict these exclusions. Each retained row has the open local
neighborhood proved in its preparation, completing the returning-arc
exhaustion.
\end{proof}

\begin{theorem}[The local secondary atlas]\label{rw:secondary-local}
The wedge and all rows of Proposition~\ref{rw:secondary:exhaustion}
have the preparations and unanchored closing-zero bounds proved above.
The bounds are two for wedge, regular, strict-real and either repeated
row, and three for finite-saddle and weak rows.
They apply with multiplicity on one actual comparison interval with
the lower whole-path and positive-central restrictions.
The adjoining maps use the exact exponent $\beta_\epsilon$;
the two lower maps use the exact common exponent $\mu$.
All coefficient contractions precede restriction to the normalized
faces, and no bound on $\epsilon|\log\lambda|$ is imposed.
\end{theorem}

\begin{proof}
Combine Propositions~\ref{rw:secondary:regular},
\ref{rw:secondary:saddle} and \ref{rw:secondary:real}.
The scalar through-domains are connected intervals, and a failed lower
whole-path condition excludes the corresponding return.
The common central maps increase in each chosen comparison input;
positive-central restrictions therefore truncate an interval and do
not add components. The exact center contributes no isolated zeros.
The returning-arc routing and tag exhaustion supply the configuration
list without using any of these counts. No conclusion here adds
the bounds from distinct chart pieces.
\end{proof}
